\scanpage{61}
\mspage{\deleted{\unsure{4|3}}}

b) Recollement de $\overline H$ à lui $=$ suivant un isom. rectifiable $[0,1]\overset{\sigma}{\longrightarrow}[-1,0]$ \inserted{$(\sigma(0)=0)$}, regardé au voisinage du pt 0.

% slanted
\margin{\uline{12.1} \deleted{\unsure{b)|a)}} c) Soit $C$ une courbe rectifiable dans \deleted{\unsure{le plan|disque}, \unsure{au bord}} \inserted{un ouvert $U$ de $\mathbb C$}, $f$ une \deleted{\unsure{germe de}} fonction complexe continue sur $U$ \deleted{\unsure{au v}} qui est continue dans $U\smallsetminus C$, alors $f$ est holomorphe dans $U$}

\textcircled{4} \uline{Recollements \deleted{bord à bord} \inserted{de disques}}. \inserted{Le recollement ``bord à bord''} Cela correspond, dans le th. 4, au cas où $K$ est une variété \inserted{(compacte)} de dim 1, $S=\varnothing$. \deleted{Il} Pour avoir ``des \unsure{nbs|ets} finis de modules'', il est naturel de dégager des conditions restrictives naturelles sur l'involution $\sigma$ \deleted{\unsure{de}} sur $\partial\widetilde X$.

Pour ceci, \deleted{rapp} on se rappelle que pour une surface conforme à bord $X$ \inserted{(au lieu de $\widetilde X$)} et une comp. connexe du bord, soit $C$, il revient au $=$ de se donner une structure de droite projective \inserted{(rectifiable)} sur $C$, ou un ``bouchage de trou'' en $C$ par un disque \inserted{de bord $C$} ; les propriétés de régularité (p.r. à la structure analytique réelle de $C$) de la structure de droite projective, reflètent
% archive sheet number: 63
